% Ad Familiares 12.18 --- headnote
% ad-familiares-12-18, late September or early October 46 BC, Rome

\textit{Cicero to Q. Cornificius, from Rome at the end of September or the
beginning of October 46 BC --- Perseus dateline \textit{Scr. Romae vel ex. m.
Sept. vel in. Oct. a. 708 (46)}. Cornificius is now in his province
(Cilicia, in the orthodox reading of this cluster), with the Caecilius
Bassus revolt smouldering in neighbouring Syria; he has written cautiously,
and Cicero approves the caution. The letter's interest is the second
section, one of Cicero's most candid sentences about life under the
dictatorship: ``civil wars always end thus, that it is not only the
victor's will that is done, but that those by whose support the victory
was gained must also be humoured.'' The names that follow --- T. Plancus
on display at Caesar's games, the mime-poets Laberius and Publilius ---
are the trivialised public culture of the new dispensation, which Cicero
endures with a face. The closing line, asking Cornificius to come back
because he is missing the one friend with whom one can laugh at these
things ``in the right intimate and learned way,'' is the warmest thing
he says in the cluster.}
